\chapter{Literature Review}
\label{ch:review}

This chapter reviews the statistical and machine-learning literature on sepsis
prediction, tightly coupled to the context and objectives of
Chapter~\ref{ch:introduction}, and formalises the three research gaps that the
thesis addresses.

\section{Statistical Onset Prediction}
TREWScore \cite{henry2015trewscore} fitted a Cox model to routinely collected
MIMIC-II vitals and laboratory values to anticipate septic shock, reaching an
AUROC of 0.83 with a median 28-hour lead time; it remains the foundational
statistical benchmark for real-time sepsis prediction. Liaw \emph{et al.}
\cite{liaw2019mortality} proposed a longitudinal, time-varying Cox extension
that tracks evolving mortality risk from repeatedly measured indicators. Kasal
\emph{et al.} \cite{kasal2004cox} compared a standard Cox model against a
flexible survival model in severe sepsis and found that hazards are frequently
\emph{non-proportional}, a methodological caution largely ignored by later
MIMIC-IV work. Critically, neither TREWScore nor its successors were fitted on
MIMIC-IV, leaving a purely statistical onset model on the current database
absent (gap \textbf{G1}).

\section{Prognostic Modelling on MIMIC-IV}
A recurring MIMIC-IV pattern couples LASSO variable selection with multivariable
logistic regression, presented as a nomogram, to predict mortality \emph{after}
diagnosis; several such studies include external validation on eICU-CRD
\cite{yuan2023nomogram}. Trajectory approaches such as group-based trajectory
modelling (GBTM) of SOFA scores \cite{yang2022gbtm} model the evolving course
but again target mortality, not onset. This sub-literature is methodologically
mature and often externally validated, but it is prognostic rather than
predictive of onset, the wrong half of the clinical problem for early warning.

\section{Benchmarks, Validation, and Reviews}
The PhysioNet/CinC 2019 Challenge \cite{reyna2020physionet} established a shared
early-prediction benchmark and a clinical-utility metric that rewards early, but
not premature, alarms. Fleuren \emph{et al.} \cite{fleuren2020ml} reviewed
machine-learning sepsis prediction and reported large between-study
heterogeneity in design and reported performance. Wang
\emph{et al.} \cite{wang2025srpm} reviewed 91 sepsis real-time prediction models
and found that the median AUROC falls from 0.886 internally to 0.783 under full
external validation, evidence that the internal-to-external gap, rather than
raw internal discrimination, is the field's central weakness (gap \textbf{G3}).
No reviewed MIMIC-IV Cox study reports testing the proportional-hazards
assumption before interpreting hazard ratios (gap \textbf{G2}).

\section{Landmarking and Informative Missingness}
Dynamic prediction by landmarking \cite{vanhouwelingen2007landmark} and its
supermodel formulation \cite{fries2023dynamiclm} pool multiple prediction
windows to emit continuously updated risk from a single interpretable model,
providing a statistical route to real-time prediction without recurrent
networks. Separately, informative missingness (the fact that \emph{whether} a
test is ordered encodes clinician suspicion) is well documented in electronic
health record modelling \cite{groenwold2020missingness,singh2021missingness}.
Both threads are central to the design of this thesis and, to the author's
knowledge, have not previously been combined for statistical Sepsis-3 onset
prediction on MIMIC-IV.

\section{Summary of Gaps}
Synthesising the above, the three gaps that structure this thesis are:
\textbf{G1}, the absence of a purely statistical real-time onset model on
MIMIC-IV; \textbf{G2}, the routine omission of proportional-hazards testing in
the MIMIC-IV Cox literature; and \textbf{G3}, the validation asymmetry whereby
external validation is common for prognostic nomograms but largely absent for
onset prediction. The contributions listed in Section~\ref{sec:contrib} address
each gap directly.
